\newglossaryentry{desidentificados}{
    name={desidentificados},
    description={Proceso mediante el cual se eliminan o transforman datos que permiten identificar a una persona de forma directa o indirecta (nombres, números de identificación, datos de ubicación, fechas) para proteger la privacidad de los individuos}
}


\newglossaryentry{esquema}{
    name={esquema},
    description={Conjunto de reglas y estructuras que definen cómo se organiza el dato en la base de datos}
}

\newglossaryentry{datos no estructurados}{
    name={datos no estructurados},
    description={ Información que carece de un formato predefinido, como fotos (.JPG), audios (MP3) o videos}
}


\newglossaryentry{normalización}{
    name={normalización},
    description={Proceso de diseño en SQL para organizar datos, eliminar duplicados y asegurar relaciones lógicas}
}


\newglossaryentry{Sharding}{
    name={Sharding},
    description={Técnica de fragmentación de datos que consiste en dividir una base de datos en partes más pequeñas (fragmentos) distribuidas entre múltiples servidores \cite{uva_nosql_sql_2019}}
}


\newglossaryentry{ACID}{
    name={ACID},
    description={Conjunto de propiedades que garantizan la fiabilidad de las transacciones en bases de datos relacionales: Atomicidad (todo o nada), Consistencia (estados válidos), Aislamiento (ejecución independiente de transacciones) y Durabilidad (persistencia de los datos ante fallos del sistema)}
}


\newglossaryentry{BASE}{
    name={BASE},
    description={Modelo que prioriza la disponibilidad y el rendimiento. El estado del sistema puede cambiar con el tiempo y la consistencia se alcanza de forma eventual}
}

\newglossaryentry{CAP}{
    name={CAP},
    description={Teorema que establece que en un sistema distribuido no es posible garantizar simultáneamente Consistencia, Disponibilidad y Tolerancia a Particiones. 
    La Consistencia implica que todos los nodos ven los mismos datos al mismo tiempo; 
    la Disponibilidad garantiza que cada solicitud recibe una respuesta, incluso en presencia de fallos; 
    y la Tolerancia a Particiones se refiere a la capacidad del sistema para continuar operando a pesar de fallos de comunicación o desconexiones entre nodos de la red. 
    Por lo tanto, se deben priorizar dos de estas propiedades según el diseño del sistema}
}

\newglossaryentry{cluster}{
    name={cluster},
    description={Conjunto de servidores interconectados que trabajan de forma coordinada para actuar como un único sistema, distribuyendo la carga de trabajo y aumentando la disponibilidad}
}

\newglossaryentry{medicamento}{
    name={medicamento},
    description={Sustancia o combinación de sustancias (fármaco + excipiente) con propiedades para prevenir, diagnosticar, tratar, aliviar o curar enfermedades en humanos o animales}
}

\newglossaryentry{vacuna}{
    name={vacuna},
    description={Sustancia o grupo de sustancias destinadas a estimular la respuesta del sistema inmunitario ante un tumor o ante microorganismos, como bacterias o virus}
}

\newglossaryentry{esavi}{
    name={evento supuestamente atribuible a la vacunación o inmunización (ESAVI)},
    description={Cualquier evento adverso que se sospecha está relacionado con la administración de una vacuna}
}

\newglossaryentry{reacción adversa a medicamentos (RAM)}{
    name={reacción adversa a medicamentos (RAM)},
    description={Según la OMS, "reacción nociva y no
                deseada que se presenta tras la administración de un medicamento, a dosis
                utilizadas habitualmente en la especie humana, para prevenir, diagnosticar o tratar
                una enfermedad, o para modificar cualquier función biológica". Nótese que esta
                definición implica una relación de causalidad entre la administración del
                medicamento y la aparición de la reacción
                }
}

\newglossaryentry{farmacovigilancia}{
    name={farmacovigilancia},
    description={Conjunto de actividades relacionadas con la detección, evaluación, comprensión y prevención de los efectos adversos o cualquier otro problema relacionado con los medicamentos}
}

\newglossaryentry{FarmaVigiC}{
    name={FarmaVigiC},
    description={Base nacional de farmacovigilancia de Cuba, gestionada por la Unidad Coordinadora Nacional de Farmacovigilancia (UCNFV) del Ministerio de Salud Pública (MINSAP)}
}

\newglossaryentry{WHO-ART}{
    name={WHO-ART (The WHO Adverse reaction terminology)},
    description={Diccionario de terminología de
                reacciones adversas de medicamentos que contiene un sistema de codificación de los
                mismos}
}

\newglossaryentry{ICD-10}{
    name={ICD-10},
    description={Es un sistema de codificación alfanumérico utilizado a nivel mundial para estandarizar la clasificación, diagnóstico, síntomas y causas externas de enfermedades, lesiones y muertes}
}

\newglossaryentry{VigiBase}{
    name={VigiBase},
    description={Base de datos internacional de vigilancia de medicamentos de la OMS}
}


\newglossaryentry{ips}{
    name={Informes Periódicos de Seguridad (IPS)},
    description={Son documentos oficiales que los laboratorios deben presentar de forma estandarizada a la autoridad reguladora (CECMED). Estos informes resumen la experiencia nacional e internacional sobre la seguridad de un medicamento en periodos definidos (cada 6 meses los primeros dos años, y luego anualmente)}
}

\newglossaryentry{fda}{
    name={FDA (Food and Drug Administration)},
    description={Agencia reguladora de medicamentos y alimentos de los Estados Unidos}
}

\newglossaryentry{formularios}{
    name={formulario},
    description={Documento estructurado que se utiliza para recopilar información de manera sistemática, siguiendo un formato predefinido}
}

\newglossaryentry{repositorios}{
    name={repositorio},
    description={Lugar centralizado donde se almacenan y mantienen datos o información de manera organizada y accesible}
}

\newglossaryentry{notificación espontánea}{
    name={notificación espontánea},
    description={Sistema de reporte voluntario en el que profesionales de la salud, pacientes o cualquier persona pueden informar sobre sospechas de reacciones adversas a medicamentos sin necesidad de una solicitud formal}
}

\newglossaryentry{faers}{
    name={FAERS (FDA Adverse Event Reporting System)},
    description={Base de datos de la FDA que recopila informes de eventos adversos relacionados con medicamentos y productos biológicos en los Estados Unidos}
}

\newglossaryentry{aems}{
    name={AEMS (Adverse Event Monitoring System)},
    description={Ecosistema inteligente desarrollado por la FDA para la monitorización de eventos adversos, que integra análisis avanzados y aprendizaje automático para detectar señales de seguridad en tiempo real}
}

\newglossaryentry{toxica}{
    name={tóxica},
    description={Relacionada con sustancias dañinas, venenosas o nocivas para la salud de los organismos vivos}
}

\newglossaryentry{biocubafarma}{
    name={BioCubaFarma},
    description={Grupo Empresarial de la Industria Biotecnológica y Farmacéutica de Cuba, que integra a todas las empresas estatales dedicadas a la investigación, desarrollo, producción y comercialización de productos biotecnológicos y farmacéuticos en el país}
}

\newglossaryentry{infranotificacion}{
    name={infranotificación},
    description={Es el fenómeno que ocurre cuando el registro de sospechas de reacciones adversas a medicamentos (RAM) es inferior a lo que sucede en la realidad.}
}

\newglossaryentry{autorreporte}{
    name={autorreporte},
    description={Proceso mediante el cual los pacientes o usuarios de medicamentos informan directamente sobre sus experiencias con reacciones adversas a medicamentos, sin intermediación de profesionales de la salud}
}

\newglossaryentry{SOLID}{
    name={SOLID},
    description={Los principios SOLID son cinco reglas de disepo de software orientado a objetos que buscan crear código más comprensible, flexible, mantenible y escalable.Estos son principio de responsabilidad única; principio de abierto/cerrado; principio de sustitución de Liskov; principio de segregación de interfaz; principio de inversión de dependencias}
}

\newglossaryentry{Repository}{
    name={Patrón Repository},
    description={Patrón de diseño que abstrae el acceso a los datos mediante una interfaz, permitiendo desacoplar la lógica de negocio de la fuente de almacenamiento}
}

\newglossaryentry{SUS}{
    name={System Usability Scale (SUS)},
    description={Cuestionario estandarizado de diez ítems utilizado para evaluar la usabilidad percibida de un sistema de forma rápida y confiable}
}

\newglossaryentry{JSON Schema}{
    name={JSON Schema},
    description={Especificación que permite definir la estructura, tipos de datos y restricciones de documentos JSON para validar su contenido}
}

\newglossaryentry{asíncrona}{
    name={asincronía},
    description={Modelo de ejecución en el que las tareas se realizan sin bloquear el flujo principal, permitiendo que múltiples operaciones se procesen de manera concurrente}
}


\newglossaryentry{vaers}{
    name={VAERS},
    description={Sistema de notificación de eventos adversos posteriores a la vacunación en Estados Unidos, utilizado para la vigilancia de la seguridad de las vacunas}
}


% \newglossaryentry{minería de datos}{
%     name={mineria de datos},
%     description={Proceso de análisis de grandes volúmenes de datos para descubrir patrones, tendencias y relaciones útiles mediante técnicas estadísticas y de aprendizaje automático}
% }


\newglossaryentry{hl7fhir}{
    name={HL7/FHIR},
    description={HL7 (Health Level Seven International) es el organismo de estandarización para el intercambio de información sanitaria. FHIR (Fast Healthcare Interoperability Resources) es un estándar moderno que define recursos y APIs REST para la interoperabilidad en salud}
}

\newglossaryentry{captcha}{
    name={CAPTCHA},
    description={Acrónimo de \textit{Completely Automated Public Turing test to tell Computers and Humans Apart}. Prueba de seguridad diseñada para distinguir entre usuarios humanos y bots automatizados, consistente típicamente en reconocer caracteres distorsionados, imágenes o resolver pequeños desafíos lógicos. En el contexto del sistema, se utiliza para prevenir el envío masivo y automatizado de reportes falsos}
}

\newglossaryentry{spam}{
    name={spam},
    description={Término utilizado para referirse a mensajes no solicitados, generalmente enviados de forma masiva y automatizada. En el contexto del sistema, se refiere al envío reiterado y malintencionado de reportes falsos o duplicados de ESAVI, lo que podría contaminar la base de datos y generar señales de alerta erróneas}
}

\newglossaryentry{rate limiting}{
    name={rate limiting},
    description={Técnica que limita el número de solicitudes por usuario o IP en un intervalo de tiempo determinado, utilizada para prevenir ataques de denegación de servicio y envío automatizado masivo de datos}
}

\newglossaryentry{rate timing}{
    name={rate timing},
    description={Técnica complementaria al rate limiting que analiza los patrones temporales entre solicitudes para detectar comportamientos automatizados (bots) que evaden los límites tradicionales}
}

\printglossaries